Sequências Numéricas

Definição: 
    Uma sequência de números reais é uma função f a valores reais cujo domínio é o conjunto dos números inteiros positivos, ou seja: f: Z+*R; n ↦ f(n) 
    Os números da imagem da sequência são chamados de elementos da sequência. 

Ex 1:
    A sequência dada pela expressão: 
        f(n)=n2n + 1
    é tal que:
        f(1) =13; f(2)=25; f(3)=37; . . . 
    ou:
        13, 25, 37, ...

Notação:
    {f(n)} : sequência determinada por f
    Denotando f(n)=an, os elementos de uma sequência também podem ser representados por:
        {an}n = 1+∞ ou {an}
    E diremos que a sequência possui um termo geral an.
    Ex 2:
    Ex 3:

Progressão aritmética (PA)
    Uma PA é uma sequência numérica {an}n = 1+∞ definida recursivamente por:
        an=an-1+r, n>1
    onde o primeiro termo a1 é um número dado e o número r é chamado de razão PA.
    Termo Geral:
        an=a1+(n-1) . r
    Tipos de PA:
PA constante: se r = 0
PA crescente: se r > 0
PA decrescente: se r < 0

Progressão Geométrica (PG)
    Uma PG é uma sequência em que cada termo, a partir do segundo, é o produto do termo anterior por uma constante q, chamada de razão da PG
        an+1=an.q
    Termo geral:
        an=a1.qn-1
    Tipos de PG:
PG constante: q = 1
PG crescente: a1>0 e q>1 ou se a1<0 e 0<q<1
PG decrescente: a1<0 e q>1 ou se a1>0 e 0<q<1
PG oscilante: q < 0

Limite de uma sequência:
    Definição
        Uma sequência {an} tem limite L e escrevemos n+∞an=L se, para cada >0 existir um inteiro correspondente N tal que se n > N, então an-L<.
        Se n+∞an existe, dizemos que a sequência {an} converge (ou é convergente). Caso contrário, dizemos que {an} diverge (ou é divergente).

    Teorema 1:
        Se x+∞f(x)=L e f(n)=an quando nZ+*, entãon+∞an=L. 
        Obs.: Se an aumentar quando n aumentar, dizemos quen+∞an= +∞. Neste caso, a sequência {an} é divergente (diverge para +∞).

    Dado o teorema 1, temos que os teoremas sobre limites de funções de uma variável real podem ser utilizados para calcular o limite de uma sequência 

    Teorema 2:
        Sejam {an} e {bn} sequências convergentes e c uma constante.
        I. n+∞(anbn)=n+∞ann+∞bn
        II. n+∞(anbn)=n+∞ann+∞bn, desde que n+∞bn0
        III. n+∞an.bn=n+∞an . n+∞bn
        IV. n+∞c . an=c .n+∞an
        V. n+∞c=c
        VI. n+∞(an)p=[n+∞an]p, p>0 e an>0

    Teorema do Confronto para funções de uma variável real.








        xI, h(x)g(x)f(x), se xaf(x)=L=xah(x), então xag(x)=L

    Teorema 3: Teorema do Confronto
        Sejam an,bn e cn sequências. Se anbncn, para nn0 e n+∞an=L=n+∞cn, então n+∞bn=L





    Teorema 4: 
        Se n+∞|an|=0, então n+∞an=0
        Ex.:
            a) an=1n
                n+∞an=n+∞1n=0
                Logo, an é convergente. Esta sequência converge para 0
            b) an=1nr
                - Se r>0: n+∞1nr=0 (Converge)
                - Se r>0: n+∞1nr=﹢∞ (Diverge)
            c) an=nn+1
                n+∞an=n+∞nn+1
                              =n+∞11+1n=1
                Logo, an converge para 1
            d) an=ln nn
                n+∞an=n+∞ln nn (Regra de L’Hôspital)
                              =n+∞1n1
                              =n+∞1n=0
                Logo, an é convergente.

    Regra de L’ Hôspital:
        I. Seja f e g funções tais que:
            xaf(x)=0 e xag(x)=0.
            Se f e g são deriváveis em a:
            xaf(x)g(x)=xaf'(x)g'(x)
        II. A regra também é válida se 
            xaf(x)= +∞ e xag(x)= +∞
            ou
            xaf(x)= -∞ e xag(x)= -∞
        Ex.:
            a) an=(﹣1)n
                n+∞an=n+∞(﹣1)n=∄, 
                Pois os valores da sequência alternam entre -1 e 1. Logo an é divergente.




             b) an=(﹣1)nn
	     Usando o teorema 4:
	     n+∞|an|=n+∞|(﹣1)nn|=0 
	     Logo, an é convergente
             c) an=sen(n.)n
                 Note que:
                 ﹣1sen(n.)1
                 0|sen(n.)|10|sen(n.)n|1n
                 Como o n+∞0=0 e n+∞1n=0, segue pelo Teorema do Confronto, que 
                 n+∞|sen(n.)n|=0, logo n+∞sen(n.)n=0 e a sequência é convergente

Teorema 5:
    Se n+∞an=L e a função f for contínua em L, então:
        n+∞f(an)=f(L) ou n+∞f(an)=f(n+∞an)
    Ex.: an=sen(n)
            n+∞sen(n)=sen(n+∞n)=sen(0)=0
            n+∞an=0 (converge)


Teorema 6:
    A sequência {rn} é convergente se -1<r1 e é divergente para todos os outros valores de r. Além disso:
    n+∞rn=0, se -1<r<1
                   =1, se r=1

Definição:
    Seja {an} uma sequência.
        i ) A sequência é crescente se an<an+1, n1
        ii ) A sequência é decrescente se an>an+1, n1
    Obs.: Uma sequência é monótona se ela for sempre crescente ou decrescente.

Definição:
    Seja uma sequência {an}:
        i ) {an} é limitada superiormente se existir um número m tal que anm, n1.
        ii ) {an} é limitada inferiormente se existir um número m tal que anm, n1.
    Obs.: Se {an} for limitada superiormente ou inferiormente, então {an} é uma sequência limitada

Teorema 7 - Sequência monótona:
    Toda sequência monótona e limitada é convergente


Séries numéricas

Definição:
    Seja {an} uma sequência. Defina:
        Sn=a1+a2+a3+ . . . +an
    Então a sequência {Sn} será chamada de série infinita, a qual é denotada por:
        n=1+∞an=a1+a2+a3+ . . . +an+ . . .
    Os números a1, a2, a3, . . ., an são os termos da série e os números S1, S2, S3,  . . . ,Sn são denominados somas parciais da sequência.
    O limite da série, quando existir, denomina-se soma da série, ou seja:
        n=1+∞an=n+∞Sn
    Se a soma for finita diremos que a série é convergente. Se a soma for infinita ou o limite não existir diremos que a série é divergente.
    
Teorema 8:
    Se a série n=1+∞an for convergente, então:
        n+∞an=0
    Obs.: 
        Deste teorema, temos o seguinte resultado:
            Se n+∞an0, então a série an é divergente

Definição - Série geométrica:
    A série geométrica é dada por:
        n=1+∞arn-1, com a= constante
    Note que n=1+∞arn-1=a+ar+ar2+ . . .
    A n-ésima soma, Sn, dessa série:
        Sn=a+ar+ar2+ . . . +arn-1
        Sn=a(1+r+r2+ . . . +rn-1) ( 1 )
    Agora, considere:
        rSn=a(r+r2+r3+ . . . +rn) ( 2 )
    Tome ( 1 ) - ( 2 )
        Sn-rSn=a(1+r+r2+ . . . +rn-1)-(a(r+r2+r3+ . . . +rn))
        Sn(1-r)=a(1-rn)
        Sn=a(1-rn)1-r




    Lembremos que uma série é convergente se o limite de Sn, quando n﹢∞, existe. Então, neste caso:
        n﹢∞Sn=n﹢∞a(1-rn)1-r
        n﹢∞Sn=n﹢∞[a1-r-a1-rrn]
        n﹢∞Sn=(a1-r-a1-r)n﹢∞rn
    Note que n﹢∞rn=0, se |r|<1
    Assim, se |r|<1, n﹢∞a1-r. 
    Portanto, se |r|<1, a série geométrica converge

. . .

    Ex.: Determine se a série n=1+∞22n31-n é convergente.
            Note que: 22n31-n=4n3-(n-1)
                                           =44n-1(13)(n-1)
                                           =4(43)(n-1)
            Assim, n=1+∞22n31-n=n=1+∞4(43)(n-1), note que a=4 e r=43. Como 43>1, então a série é divergente 

Teorema 9
    Sejam n=1+∞an e n=1+∞bn séries convergentes e c uma constante. Então:
        i ) n=1+∞can é convergente e n=1+∞can=cn=1+∞an.
        ii ) n=1+∞an n=1+∞bn é convergente
        Ex.: Determine se a série n=1+∞(3n(n+1)+12n) é convergente.
            Note que n+∞(3n(n+1)+12n)=0, logo é inconclusivo pelo teorema 7.
            Mas: n=1+∞(3n(n+1))=3n=1+∞(1n(n+1)). Logo, a série n=1+∞(1n(n+1)) é convergente (Como visto no exemplo da aula passada)
            Além disso:
                12n=(12)n=(12)(12)n-1
            Assim:
                n=1+∞12n=n=1+∞(12)(12)n-1
            É uma série geométrica com r=12<1, convergente.
            Portanto, n=1+∞(3n(n+1)+12n)=n=1+∞(1n(n+1))+n=1+∞(12)(12)n-1 é uma série convergente

Teorema 10
    Sejam n=1+∞an uma série divergente, e n=1+∞bn uma série convergente e c uma constante, c0. Então:
        i ) n=1+∞can é divergente.
        ii ) n=1+∞(an+bn)  é divergente.

Teorema 11: Testes de comparação
    Suponha que an e bn sejam séries de termos positivos.
    1. Se bn for convergente e anbnn, então an também será convergente.
    2. Se bn for divergente e anbnn, então an também será divergente.
    Ex.: Determine se a série é convergente ou divergente.
        a) n=1+∞43n+1
            Note que, para n1:
                3n<3n+113n>13n+143n>43n+1
            Além disso: 
                43n=4313n-1=43(13)n-1


            Além disso, temos a série geométrica: 
                n=1+∞43(13)n-1, que é convergente pois r=13<1
            Como 43n+1<43n=43(13)n-1, para todo n1, segue que a série n=1+∞43n+1 é convergente (pelo teste da comparação)

        b) n=1+∞1n
            Para n1:
                nn1n1n
            Como a série harmônica n=1+∞1n é divergente, pelo teste da comparação temos que a série n=1+∞1n, também é divergente.

Definição: Série p (Série hiper-harmônica)
    A série:
        n=1+∞1np
    Teorema 12:
        A série p:
            1. Converge, se p>1
            2. Diverge, se p1
        Demonstração:
            1. Se p=1, temos a série harmônica, que é divergente
            2. Se p=0: n=1+∞1np=n=1+∞1, n+∞1=10. Logo esta série é divergente.
            3. Se p<0: n=1+∞1np, n+∞1np=﹢∞, pois p<0. Esta série diverge também.
            4. Se p>1: n=1+∞1np=11p+12p+13p+14p+15p+16p+17p+18p+ . . .
                Note se que se n1:
                    - 2p<3p12p>13p
                    - 4p<5p14p>15p
                Assim:
                    n=1+∞1np<11p+12p+12p+14p+14p+14p+14p+18p+ . . .
                    n=1+∞1np<11p+22p+44p+88p+ . . . n=1+∞1np<20(20)p+21(21)p+22(22)p+23(23)p+ . . . 
                    n=1+∞1np<n=1+∞2n-1(2n-1)p=n=1+∞2n-1(2p)n-1 
                    n=1+∞1np<n=1+∞(22p)n-1=n=1+∞(12p-1)n-1 

                Temos que:
                    n=1+∞(12p-1)n-1 é uma série geométrica com a=1 e r=12p-1, se p>1, p-1>0, então r=12p-1<1 e esta série é convergente. 
                Pelo teste de comparação, como n=1+∞1np<n=1+∞(12p-1)n-1, temos que a série p é convergente para n>1

Teorema 13 - Teste de comparação de limite:
    Suponha que an e bn sejam séries de termos positivos.
        i ) Se n+∞anbn=c>0, então ambas as séries são convergentes ou divergentes.
        ii ) Se n+∞anbn=0 e se bn converge, então an também converge.
        iii ) Se n+∞anbn=﹢∞ e se bn diverge, então an também diverge.

Séries alternadas
    n=1+∞an é uma série alternada se an=(﹣1)nbn ou an=(﹣1)n-1bn, onde bn é um número positivo (|an|=bn)



Teorema 14 - Teste da série alternada
    Se a série é alternada n=1+∞(﹣1)n-1bn=b1-b2+b- . . . , com bn>0, satisfaz:
        i ) bn+1<bn, n (decrescente)
        ii ) n+∞bn=0
    Então a série é convergente.

Convergência absoluta
    Seja an, uma série qualquer. Temos que:
        |an|=|a1|+|a2|+ . . .
    É a série cujos termos são os valores absolutos da série original.

    Definição:
        Uma série an é dita absolutamente convergente se a série de valores absolutos |an| for convergente.
    Definição:
         Uma série é denominada condicionalmente convergente se ela for convergente, mas não absolutamente convergente.

Teorema 15:
    Se a série for absolutamente convergente, então an é convergente
Teorema 16 - Teste da razão:        (Mais recomendado quando tem algum fatorial)
    Seja a série an e considere que exista o limite:
        n+∞|an+1an|=L
    1. Se L<1, a série é absolutamente convergente, e portanto é convergente.
    2. Se L>1, a série é divergente.
    3. Se L=1, o teste é inconclusivo
    Obs.: Se n+∞|an+1an|=﹢∞, então a série an é divergente.
    Ex.: Use o teste da razão para determinar se a série dada é convergente ou divergente
            a) n=1+∞(﹣1)nn33n
                 (Resolução na foto do dia 17/09)
                 . . . Portanto, a série dada é absolutamente convergente e, consequentemente, convergente
            b) n=1+∞1n
                (Resolução na foto do dia 17/09)
                . . . Neste caso o teste é inconclusivo (Sabemos que esta série é divergente por exemplos anteriores)
            c) n=1+∞1n2
                (Resolução na foto do dia 17/09)
                . . . Neste caso o teste é inconclusivo (Sabemos que esta série é convergente por ser uma série p com p>1)
            d) n=1+∞nnn!
                (Resolução na foto do dia 17/09)
                . . . e>1, portanto, esta série é divergente

Teorema 17 (Teste da Raiz):                  (Mais recomendado quando toda a expressão está elevada a n)
    Seja an uma série e considere que exista o limite:
        n+∞n|an|=L
    1. Se L<1, a série é absolutamente convergente e, portanto, convergente.
    2. Se L>1, a série é divergente.
    3. Se L=1, o teste é inconclusivo
    Obs.: Se n+∞n|an| não existir, então a série é divergente.
    Ex.: Use o teste da raiz para determinar se a série é convergente ou divergente.
            n=1+∞(2n + 33n + 2)n
            (Resolução foto do dia 17/09)
            . . . Pelo teste da raiz, a série dada é convergente


Apêndice -> Integral imprópria:
    . . . Gráfico
    A(t)=1t1x2dx=﹣1x|x=1x=t=﹣1t-(﹣11)
    A(t)=1-1t      A(t)<1
    t+ ∞A(t)=t+ ∞(1-1t)=1
   A área abaixo da curva de f(x)=1x2, para x>1, se aproxima de 1 quando t﹢∞. Assim:
        1+∞1x2 dx=t+∞1t1x2 dx
    Definição:
        i ) Se atf(x) dx existe para cada número ta, então
                a+∞f(x) dx=t+∞atf(x) dx, desde que este limite exista

        ii ) Se tbf(x) dx existe para cada tb, então:
                -∞bf(x) dx=t-∞tbf(x) dx, desde que este limite exista
        As integrais impróprias a+∞f(x) dx e -∞bf(x) dx são chamadas de convergentes se os correspondentes existirem e divergentes se não existirem.


        iii ) Se ambas as integrais a+∞f(x) dx e tbf(x) dx forem convergentes, então definimos:
                -∞+∞f(x) dx=-∞af(x) dx+a+∞f(x) dx

Teorema 18 - Teste da Integral: 
    Suponha que f seja uma função contínua positiva e decrescente em [1, ﹢∞[ e seja an=f(n). Então:
    1) Se 1+∞f(x) dx for convergente, então n=1+∞an é convergente
    2) Se 1+∞f(x) dx for divergente, então n=1+∞an é divergente
    Obs.:
        i ) Se n=k+∞an, k>1, então teremos a integral k+∞f(x) dx.
        ii ) Não é necessário que f seja sempre decrescente, o importante é que ela seja a partir de algum N>1. Se n=N+∞an é convergente, então n=1+∞an é convergente. Logo é possível aplicar o teste da integral
    Ex.: Usando teste da integral
        a) n=1+∞1n2+1
            Seja f(x)=1x2+1
            Note que f é contínua e positiva, x1. Além disso, para x1: 
                x2<x2+11x2>1x2+1
            Logo, f é decrescente.
            Outra forma de mostrar que f é decrescente:
                f(x)=1x2+1=(x2+1)-1
                f'(x)=﹣(x2+1)-2(2x)
                f'(x)=﹣2x(x2+1)2
            Como x1 (positivo), então ﹣2x é negativo x1. Como (x2+1)2 é sempre positivo, que 
            f'(x)<0 (f é decrescente), para x1. Assim, o teste da integral é aplicável.
                1+∞1x2+1 dx=t+∞1t1x2+1 dx
                                 =t+∞[tg-1 (x)| x=1x=t ]
                                 =t+∞[tg-1(t)-tg-1(1)]
                                 =t+∞[tg-1(t)-4]=t+∞tg-1(t)-t+∞4

                                 =2-4=4
         Logo, 1+∞1x2+1 dx=4, a integral imprópria é convergente.
         Portanto, a série dada é convergente.


        b) n=1+∞ln nn
            Seja f(x)=ln xx. Para x1, f é contínua e positiva.
                f'(x)=(ln x)'  x - (ln x)  x'x2
                f'(x)=1x  x - (ln x)  1x2=1 - ln xx2
                f'(x)<01-ln x <01<ln x
                e1<eln xe<xx>e
            Logo, f é decrescente para x>e. Assim, pela observação ii, é possível aplicar o Teste da Integral
                1+∞ln xx dx=t+∞1tln xx dx
             Vamos calcular a integral I
                 I=1tln xx dx    { du = 1x dxu = ln x
                 I*=u du=u22+C1
                 I=[ln x]22 |  x = 1x = t=[ln t]22-[ln 1]22
             Assim:
                 I=1tln xx dx=[ln t]22
             Logo:
                 1+∞ln xx dx=t+∞[ln t]22=﹢∞
             Como a integral imprópria é divergente, segue que a série dada diverge.

















































































































Limite
Definição:
    Seja f uma função de duas variáveis que está definida em uma bola aberta B((a,b),r), exceto possivelmente em (a,b). Então o limite de f(x,y) quando (x,y) tendem a (a,b) é L, e escrevemos 
            (x,y)(a,b)f(x,y)=L
    Se >0, >0; 0<||(x,y)-(a,b)||<, então |f(x,y)-L|<, 
    onde ||(x,y)-(a,b)||=(x-a)2+(y-b)2
    Outras notações para limite:
        yxbaf(x,y)=L com duas setinhas (xa e yb)
      ou
        f(x,y)L quando (x,y)(a,b)


            Gráfico



Teorema:
    Se f(x,y)=L1 quando (x,y)(a,b) ao longo do caminho C1 e f(x,y)=L2 quando (x,y)(a,b) ao longo do caminho C2, com L1L2, então (x,y)(a,b)f(x,y) ∄ 
    ex.: Mostre que
        a) (x,y)(0,0)x2 - y2x2 + y2 não existe
            Caminho C1: y=0
                (x,y)(0,0)x2 - 0x2 + 0=x2x2=1
            Caminho C2: x=0
                (x,y)(0,0)0 - y20 + y2=-y2y2=﹣1

            Logo, o limite não existe
    







Os teoremas sobre limites de funções de uma variável podem ser estendidos para as funções de duas variáveis. Em particular:
    (x,y)(a,b)x=a
    (x,y)(a,b)y=b
    (x,y)(a,b)c=c, com c=constante

Teorema do Confronto:
    Se f(x,y)g(x,y)h(x,y), (x,y)B((a,b),r) e se 
            (x,y)(a,b)f(x,y)=(x,y)(a,b)h(x,y)=L
    então:
            (x,y)(a,b)g(x,y)=L

Teorema - Função Limitada:             ↱g é uma função limitada
    Se (x,y)(a,b)f(x,y)=0 e se |g(x,y)|<M, (x,y)B((a,b),r), onde r>0 e M>0 são números reais fixos, então
            (x,y)(a,b)f(x,y)g(x,y)=0

Teorema:
    (x,y)(a,b)f(x,y)=0(x,y)(a,b)|f(x,y)|=0  

   exercícios

Continuidade:
    Dizemos que uma função f de duas variáveis x e y é contínua no ponto (a,b) se, e somente se, as seguintes condições forem satisfeitas:
        i ) f(a,b) , ou seja, (a,b)Dm(f)
        ii ) (x,y)(a,b)f(x,y) existe
        iii ) (x,y)(a,b)f(x,y)=f(a,b)
    Se pelo menos uma dessas não for satisfeita, então dizemos que f é descontínua em (a,b)
    Obs.: Se uma função f for descontínua em (a,b), mas (x,y)(a,b)f(x,y) , então f possui uma 
              descontinuidade removível em (a,b), pois se redefinirmos f como segue:
              f1(x,y)={ (x,y)(a,b)f(x,y) se (x,y) = (a,b)f(x,y), se (x,y)  (a,b), então f1 será contínua em (a,b). 

              Caso contrário, dizemos que f possui uma descontinuidade essencial em (a,b).
    ex.: f(x,y)=x3x2 + y2
            Note que (0,0)Dm(f). Assim, f(0,0) não existe, e portanto, f é descontínua em (0,0). Mas, em 
            um exemplo anterior, vimos que:
            (x,y)(a,b)x3x2 + y2=0
            Neste caso, a descontinuidade é removível
            f1(x,y)={ 0 se (x,y) = (0,0)x3x2 + y2, se (x,y)  (0,0)

Derivada parcial
Definição
    Se f é uma função de duas variáveis x e y, suas derivadas parciais são as funções fx e fy, definidas por:
        i ) fx(x,y)=h0f(x+h, y) - f(x, y)h
        ii )fy(x,y)=h0f(x, y+h) - f(x, y)h
Notação para derivadas parciais: se z=f(x,y)
    i ) fx(x,y)=fx=𝜕f𝜕x=𝜕𝜕xf(x,y)=𝜕z𝜕x=f1=D1f=Dxf
    ii ) fy(x,y)=fy=𝜕f𝜕y=𝜕𝜕yf(x,y)=𝜕z𝜕y=f2=D2f=Dyf

ex.:











Derivadas de ordem superior
    Se z=f(x,y) as derivadas parciais de 2a ordem de f são as derivadas parciais de fx e fy.

    Usamos a seguinte notação:
        i ) (fx)x=fxx=𝜕𝜕x(𝜕f𝜕x)=𝜕2f𝜕x2=𝜕2z𝜕x2
        ii ) (fy)y=fyy=𝜕𝜕y(𝜕f𝜕y)=𝜕2f𝜕y2=𝜕2z𝜕y2
        iii ) (fx)y=fxy=𝜕𝜕y(𝜕f𝜕x)=𝜕2f𝜕y𝜕x=𝜕2z𝜕y𝜕x
        iv ) (fy)x=fyx=𝜕𝜕x(𝜕f𝜕y)=𝜕2f𝜕x𝜕y=𝜕2z𝜕x𝜕y

Teorema de Clairaut:
    Suponha que f seja definida em uma bola aberta D que contenha o ponto (a,b). Se as funções fxy e fyx forem ambas contínuas em D, então:
        fxy(a,b)=fyx(a,b)

    ex.:

    Derivadas de 3a ordem também podem ser definidas. 
    Por exemplo:
        fxxy=𝜕𝜕y(𝜕2f𝜕x2)=𝜕3f𝜕y𝜕x2
    O Teorema de Clairaut é aplicável a derivadas parciais de ordem maior ou igual a 3
    ex.: Seja a função . . .


Plano Tangente
Definição:
    Suponha que f tenha derivadas parciais de 1a ordem contínuas. Uma equação do plano tangente a superfície z=f(x,y), no ponto P(a,b,f(a,b)) é dada por:
        z=f(a,b)+fx(a,b)(x-a)+fy(a,b)(y-b)
    ex.: Determine a equação do plano tangente a superfície da função f(x,y)=4-x2-2y2, no ponto P(1,1,1) . . . 


Aproximação Linear
A função linear L1 dada por: 
        L(x,y)=f(a,b)+fx(a,b)(x-a)+fy(a,b)(y-b)
    é denominada linearização de f em (a,b) e a aproximação:
        f(x,y)L(x,y)
    é chamada de aproximação linear ou aproximação pelo plano tangente de f em (a,b).
    ex.: Determine a linearização L de f em (1,4), onde f(x,y)=yln x

questão da prova: especificar o domínio e fazer o esboço

   . . . . . .


Diferenciais
Seja f uma função de duas variáveis, tal que z=f(x,y). Suponha que x varie de a para a+x e y varie de b para b+y. Então, o incremento de z é z, dado por:
    z=f(a+x, b+y)-f(a,b)

Definição:
    Se z=f(x,y), então f é diferenciável em (a,b) se z pode ser expresso na forma:
        z=fx(a,b)x+fy(a,b)y+1x+2y
    onde 10 e 20 quando (x,y)(0,0)
    Obs.: Esta definição afirma que uma função diferenciável é aquela para a qual a aproximação linear é uma 
              boa aproximação quando (x,y) está próximo de (a,b)



Teorema 1:
    Seja f uma função de duas variáveis. Se f for diferenciável em (a,b) então f é uma função contínua em (a,b)
    Sejam f:AR2R, A uma bola aberta e (a,b)A

Teorema 2:
    Se f for diferenciável em (a,b), então f admitirá derivadas parciais nesse ponto.

Teorema 3:
    Se as derivadas parciais fx e fy existirem em A e forem contínuas em (a,b), então f será diferenciável em (a,b)

Obs.:
    i ) Se f não for contínua em (a,b), então f não será diferenciável em (a,b)
    ii ) Se fx ou fy não existir em (a,b), então f não será diferenciável em (a,b)

ex.: Seja f(x,y)=xexy
















Seja f uma função de duas variáveis z=f(x,y). Definimos as diferenciais dx e dy como variáveis independentes e a diferencial dz, chamada de diferencial total, dada por:
    dz=fx(x,y)dx+fy(x,y)dy
Notação: dz=df
Se tomarmos:
    dx=x=x-a
    dy=y=y-b
Então:
    dz=fx(a,b)(x-a)+fy(a,b)(y-b)
Assim, com a notação de diferencial, a aproximação linear pode ser escrita como:
    f(x,y)f(a,b)+dz

ex.:






Derivada direcional e o vetor gradiente
Definição:
    A derivada direcional de f em (x0,y0) na direção do vetor unitário u=(a,b) é:
            Duf(x0,y0)=h0f(x0 + ha, yo + hb) - f(x0, y0)h
    se este limite existir

Notação:
    Duf(x0,y0)= ∂f∂u(x0,y0)
    obs.: 
        i ) A derivada direcional de f na direção de u no ponto (x0,y0) é o coeficiente angular da reta tangente à 
            curva C no ponto (a,b,f(a,b)) do gráfico f
        ii ) As derivadas parciais fx e fy são casos particulares da derivadas direcionais quando u=e1=(1,0) 
            e u=e2=(0,1)
ex.: Seja f(x,y)=3x2-y2+4x e u o vetor unitário direcionado na direção 6. Determine Duf(x0,y0), com (x0,y0)ℝ

…

Definição:
    Se f é uma função diferenciável de x e y, o gradiente de f é a função vetorial ∇f, definida por:
            ∇f(x,y)=(fx(x,y),fy(x,y)) ou seja, ∇f(x,y)=fx(x,y)i+fy(x,y)j

Teorema:
    Se f é uma função diferenciável em x e y, então f possui derivada direcional na direção de qualquer vetor u=(a,b) e:
        Duf(x,y)=∇f(x,y)u

ex.:












Teorema:
    Suponha que f seja uma função de 2 ou 3 variáveis e xDom(f). O valor máximo da derivada direcional Duf(x,y) é ||∇f(x,y)|| e ocorre quando o vetor u tem a mesma direção do vetor gradiente ∇f(x)
    Demonstração:
        Duf(x)=∇f(x)u
                   =||f(x)||||u|| . cos()
        Como 0 (entre vetores), o maior valor de cos()=1, quando =0 e se =0, ∇f(x) e u têm as mesma direção e Duf(x)=||f(x)|| é o valor máximo da derivada direcional

Regra da Cadeia (Seção 14.5)
Para funções de uma variável real, se y=f(u) e u=g(x), então:
    dydx=dydududx
ou
    Dx[f(g(x))]=f'(g(x))g'(x)

Para funções com mais de uma variável real há várias versões da Regra da Cadeia.

Caso 1:
    Seja z=f(x,y) uma função diferenciável de x e y, onde x=g(t) e y=h(t) são funções diferenciáveis de t. Então z é uma função diferenciável de t e:
        dzdt=∂z∂xdxdt+∂z∂ydydt

    Ex.: Se z=x2y+3xy4, onde x=sen(2t) e y=cos(t), determine dzdt quando t=0
            dzdt=∂z∂xdxdt+∂z∂ydydt
                =((2xy+3y4)2(cos(2t)))+((x2+12xy3))﹣sen(t))
                =(4xy+6y4)cos(2t)-(x2+12xy3))sen(t)
            Quando t=0:
                cos(20)=cos(0)=1
                sen(0)=0
                dzdt=(4xy+6y4)cos(2t)-(x2+12xy3))sen(t)
                    =(4xy+6y4)1-(x2+12xy3))0
                    =4xy+6y4      x=sen(20)=0;y=cos(0)=1
                    =401+614=6
            Logo, quando t=0, dzdt=6

Caso 2:
    Seja z=f(x,y) uma função diferenciável de x e y, onde x=g(s,t) e y=h(s,t) são funções diferenciáveis de s e t. Então z é uma função diferenciável de s e t e:
        i ) ∂z∂s=∂z∂x∂x∂s+∂z∂y∂y∂s
        ii ) ∂z∂t=∂z∂x∂x∂t+∂z∂y∂y∂t
    Neste caso, há 3 tipos de variáveis: s e t são as variáveis independentes, x e y são as variáveis intermediárias e z é a variável dependente.



Caso geral:
    Suponha que u seja uma função diferenciável de n variáveis x1,x2,x3, . . . ,xn, onde cada xj é uma função diferenciável de m variáveis t1,t2,t3, . . . ,tm. Então u é uma função diferenciável de t1,t2,t3, . . . ,tm e:
        ∂u∂ti=∂u∂x1∂x1∂ti+∂u∂x2∂x2∂ti+∂u∂x3∂x3∂ti+ . . .+∂u∂xn∂xn∂ti, para i=1,2,3, . . . ,m
    Ex.: Se u=x4y+y2z3, onde x=rset,y=rs2e-t e z=r2ssen(t), determine o valor de
            ∂u∂s|r = 2, s = 1, t = 0
            ∂u∂s=∂u∂x∂x∂s+∂u∂y∂y∂s+∂u∂z∂z∂s
                =(4x3y)(ret)+(x4+2yz3)(2sre-t)+(3y2z2)(r2sen(t))
            Para r=2,s=1 e t=0
            ∂u∂s=4x3y2e0+(x4+2yz3)22e-0+3y2z222sen(0)
                =8x3y+4x4+8yz3
            x=21e0=2;y=212e-0=2; e z=221sen(0)=0
            ∂u∂s=8232+424+8203
                =128+64
            Logo, ∂u∂s|r = 2, s = 1, t = 0=192

Valores máximo e mínimo (Seção 14.7)
Definição:
    Uma função de duas variáveis tem máximo local (relativo) em (a,b) se f(x,y)f(a,b), (x,y) em uma bola aberta com centro (a,b), ou seja, quando (x,y) está próximo de (a,b). O número f(a,b) é chamado máximo local.
    Se f(x,y)f(a,b) quando (x,y) está próximo de (a,b), então f tem um mínimo local em (a,b) e f(a,b) é um valor mínimo local.
    Obs.: Se as inequações da definição anterior valerem (x,y)Dm(f) , então f tem máximo global (absoluto) ou tem máximo global (absoluto) em (a,b)

Teorema:
    Se f tem um máximo ou mínimo local em (a,b) e as derivadas parciais de 1a ordem de f existirem em (a,b) então ambas são nulas, ou seja fx(a,b)=0 e fy(a,b)=0
Obs.: A interpretação geométrica do teorema acima é que o gráfico de  f tem um plano tangente horizontal em um máximo ou mínimo local

    




Definição:
    Um ponto (a,b) é chamado ponto crítico de f se fx(a,b)=0 e fy(a,b)=0, ou se uma das derivadas de 1a ordem não existir
    Ex.: f(x,y)=4-x2-y2
            Derivadas parciais de f:
                fx(x,y)=﹣2x
                fy(x,y)=﹣2y
            {fy(x,y) = 0fx(x,y) = 0{-2y = 0-2x = 0{y = 0x = 0
            Logo, há somente um ponto crítico de f em (0,0)
        Note que:
            z=f(x,y)z=4-(x2+y2) e x2+y20
        Assim, f assume valor máximo absoluto quando x2+y2=0, ou seja, em x=0 e y=0

Teste da derivada segunda:
    Seja f uma função de duas variáveis x e y e supunha que as suas derivadas parciais de 2a ordem sejam contínuas em uma bola aberta de centro em (a,b). Além disso, suponha que fx(a,b)=0 e fy(a,b)=0, ou seja, que (a,b) é um ponto crítico de f. Seja
        D=D(a,b)=fxx(a,b)fyy(a,b)-[fxy(a,b)]2
    i ) Se D>0 e fxx(a,b)>0, então f(a,b) é um mínimo relativo
    ii ) Se D>0 e fxx(a,b)<0, então f(a,b) é um máximo relativo6cx
    iii ) Se D<0, então f(a,b) não é mínimo relativo nem máximo relativo. Neste caso, (a,b) é denominado 
          ponto de sela de f 
    iv ) Se D=0, o teste é inconclusivo.

    Obs.: Podemos escrever D como um determinante
                D=D(a,b)= |fyxfxx   fyyfxy| =fxxfyy-[fxy(a,b)]2
    Lembrando que:
        fxy=fyx fxy e fyx forem contínuas

    Ex.: 








Teorema do valor extremo para funções de duas variáveis:
    Seja f uma função de duas variáveis x e y, contínua no conjunto fechado Aℝ2. Então existem (x1,y1) e (x2,y2) em A tais que:
        f(x1,y1)f(x,y)f(x2,y2),(x,y)A
    Obs.:
        Este teorema afirma que f possui um mínimo absoluto em (x1,y1) e um máximo absoluto em (x2,y2), no 
        conjunto fechado A.
    Para determinar os extremos absolutos de uma função contínua f em um conjunto fechado A, podemos seguir os passos abaixo:
        1) Determine os valores de f nos pontos críticos de f em A.
        2) Determine os valores externos de f na fronteira de A.
        3) O maior valor dos passos 1 e 2 é o máximo absoluto de f, e o menor, é o mínimo absoluto.

    Ex.: 






































Equação Diferencial Ordinária de 1ᵃ Ordem (Guidorrizzi, Volume 4)
Uma equação diferencial ordinária (EDO) de 1ᵃ ordem é uma equação do tipo:
    dydx=F(x,y)       (1)
Onde F é uma função definida em um aberto Dℝ2.
Dizemos que y=y(x), definida em um aberto Iℝ2, é uma redução da equação (1) se, xI, ocorrer:
    y'(x)=F(x,y(x))

Ex.: Verifique que y(x)=ex2, xℝ, é uma solução da EDO:
          dydx=2xy

EDO de 1ᵃ ordem de variáveis separáveis
    Esse tipo de equação tem a forma:
        dydx=g(x)h(y)       (2)
    Onde g e h são funções definidas em intervalos abertos I1 e I2 respectivamente.
    Ex.:
        a) dydx=2xy2 → Separável: g(x)=x e y(x)=y2
        b) dydx=x2+y2 → Não é separável
    Soluções constantes de uma EDO de 1a ordem de variáveis separáveis
        Suponha que g(x)0,xI1. A função constante y(x)=K, com xI2, será uma solução da equação (2) se K for raiz da equação h(y)=0
        Ex.:

        Ex.:

    Soluções não constantes de uma EDO de 1a ordem de variáveis separáveis
        dydx=g(x)h(x)1h(x)dy=g(x)dx1h(x)dy=g(x)dx+c, onde c é uma constante.
        Ex.:

        Ex.: PVC
